% Copyright (c) 2026 Caio Terencio. All rights reserved.
% 
% Licensed under the Apache License, Version 2.0 (the "License");
% you may not use this file except in compliance with the License.
% You may obtain a copy of the License at
%
%     http://www.apache.org/licenses/LICENSE-2.0
%

% Início do documento LaTeX: define tipo e formato do documento
% a4paper = tamanho A4, 10pt = tamanho base da fonte
\documentclass[a4paper,10pt]{article}

% --- PACOTES: BASE / IDIOMA ---
\usepackage[utf8]{inputenc}
\usepackage[T1]{fontenc}
\usepackage[portuguese]{babel}

% --- PACOTES: LAYOUT / TIPOGRAFIA ---
\usepackage{geometry}
\usepackage{parskip}
\usepackage{microtype}
\usepackage{enumitem}
\usepackage{titlesec}
% Os pacotes array e tabularx foram removidos para evitar problemas de parsing em ATS

% --- PACOTES: LINKS / BOOKMARKS ---
\usepackage{hyperref}
\usepackage{bookmark}
\usepackage{xurl}

% --- CONFIGURAÇÃO: MARGENS / ESPAÇAMENTO ---
\geometry{top=1.5cm, bottom=1.5cm, left=1.5cm, right=1.5cm} % Margens da página
\setcounter{secnumdepth}{0}                                 % Remove numeração de seção
\setlist[itemize]{
    leftmargin=0.75em, % Recuo da lista (mais perto da margem da página)
    itemsep=0.2em,     % Espaço entre itens
    topsep=0.25em,     % Espaço antes/depois da lista
    parsep=0em,        % Espaço entre parágrafos no mesmo item
    partopsep=0em      % Espaço extra ao iniciar a lista
}

% Remove números de página e cabeçalhos para visual limpo do currículo
\pagestyle{empty}

% --- CONFIGURAÇÃO: METADADOS / LINKS ---
\hypersetup{
    pdftitle={CV Caio Terencio}, % Título do PDF
    pdfauthor={Caio Terencio Trindade Ribeiro}, % Autor do PDF
    colorlinks=true,          % Usa cor nos links (sem caixa)
    linkcolor=black,          % Cor de links internos
    urlcolor=black,           % Cor de URLs
    citecolor=black,          % Cor de citações
    bookmarksdepth=2          % Barra lateral: seção e subseção
}

% --- CONFIGURAÇÃO: TÍTULOS ---
\titleformat{\section}
{\Large\bfseries}
{}
{0em}
{}
[\titlerule\vspace{0.5ex}]
\titleformat{\subsection}
{\normalsize\bfseries}
{}
{0em}
{}
\titlespacing*{\subsection}{0pt}{0.35em}{0.15em}

% Macro para entradas do currículo com bookmark no PDF (Otimizado para ATS - Layout Linear)
\newcounter{cventry}
\newcommand{\cventry}[4]{%
    \refstepcounter{cventry}%
    \phantomsection%
    \pdfbookmark[2]{#1}{cventry-\thecventry}%
    {\noindent\textbf{#1} \hfill #2\par}
    {\noindent\textit{#3} \hfill \textit{#4}\par}
    \vspace{0.2em}
}

% --- INÍCIO DO DOCUMENTO ---
\begin{document}

% --- CABEÇALHO ---
\begin{center}
    {\LARGE \textbf{Caio Terencio Trindade Ribeiro}} 
    \\ [0.1cm]
    São Paulo, SP
    {\textbullet}
    \href{mailto:caioterencio@gmail.com}{caioterencio@gmail.com} 
    {\textbullet}
    \href{https://linkedin.com/in/caioterencio}{linkedin.com/in/caioterencio} 
    {\textbullet}
    \href{https://github.com/caio0708}{github.com/caio0708}
\end{center}

% --- RESUMO PROFISSIONAL ---
\section{Resumo Profissional}
Engenheiro de Controle e Automação com mais de 2 anos de experiência prática no desenvolvimento de soluções integrando software, hardware, IoT e inteligência artificial Experiência com hiperautomação utilizando n8n, desenvolvimento de agentes de IA (LLMs), APIs REST e back-end em Python. Também possuo experiência com ESP32, sensores, eletrônica e soluções IoT além de RPA e Machine Learning aplicados à otimização de processos e operações.



% --- SEÇÕES ---

\section{Experiência Profissional}
    
    \cventry{ABMIX}{São Paulo, SP (Remoto)}{Analista de Automações}{Dez 2025 - Presente}
        \begin{itemize}
            \item Arquitetura e gestão de hiperautomações utilizando n8n, implementando tratamento robusto de dados para evitar falhas de execução por parâmetros nulos em integrações sistêmicas.
            \item Desenvolvimento de agentes de IA e chatbots avançados, empregando LLMs e máquinas de estado para roteamento inteligente de requisições, minimizando escalonamentos desnecessários para o atendimento humano.
            \item Criação de fluxos automatizados para captação, qualificação e nutrição de leads, conectando plataformas diversas através do consumo de APIs REST.
            \item Estruturação de ambientes de gestão no ClickUp, desenvolvendo dashboards dinâmicos e painéis de monitoramento para análise de leads qualificados e churn de campanhas.
            \item Implementação de soluções de RPA (Robotic Process Automation) para extração automatizada de dados e boletos comerciais.
        \end{itemize}

    \cventry{Siemens Healthineers}{São Paulo, SP (Híbrido)}{Estagiário - Terapias Avançadas}{Mai 2024 - Nov 2025}
        \begin{itemize}
            \item Programação e agendamento de equipamentos de demonstração.
            \item Suporte na configuração e precificação de produtos e apoio a apresentações.
            \item Auxílio na elaboração de propostas técnicas para licitações e cadastro de itens no SAP.
            \item Auxílio em atividades de inteligência de mercado e análise de relatórios da concorrência.
            \item Suporte em projetos especiais, feiras, congressos e desenvolvimento de materiais de marketing.
        \end{itemize}

% ------

\section{Educação}
    
    \cventry{FIAP}{São Paulo, SP}{Pós Graduação em Machine Learning Engineering}{Jan 2026 - 2027}
    
    \cventry{Instituto Mauá de Tecnologia}{São Caetano do Sul, SP}{Bacharelado em Engenharia de Controle e Automação}{Jan 2021 - Dez 2025}

% ------

\section{Experiência Acadêmica e Projetos}

    \cventry{Instituto Mauá de Tecnologia}{São Caetano do Sul, SP}{Iniciação Científica - Pesquisador}{Jan 2023 - Dez 2023}
        \begin{itemize}
            \item Desenvolvi um aplicativo móvel em Python utilizando Kivy para monitoramento de telemetria em tempo real de embalagens logísticas de transplantes.
            \item Programei firmware de ESP32 integrando aquisição de dados via MQTT e plataforma Node-RED.
            \item Vencedor do prêmio ABINC de IoT na categoria acadêmica durante a conferência Futurecom 2024.
        \end{itemize}

    \cventry{GCSP (Grand Challenges Scholars Program) - IMT}{São Caetano do Sul, SP}{Grand Challenges Scholar - Diamond}{Jan 2023 - Dez 2025}
        \begin{itemize}
            \item Recebi a Carta de Comenda da National Academy of Engineering (NAE), dos Estados Unidos, em reconhecimento à conclusão do nível Diamond do GCSP e ao desenvolvimento das cinco competências do programa: técnico-criativa, multidisciplinar, empreendedorismo, multicultural e consciência social.
        \end{itemize}

    \cventry{Projetos Independentes e Pessoais}{São Paulo, SP}{Desenvolvedor Full-Stack e Engenheiro IoT}{Jan 2026 - Presente}
    \begin{itemize}
        \item Desenvolvimento de sites e aplicações full-stack (React, TypeScript e back-end), integrando APIs REST para plataformas interativas e sistemas de monitoramento contínuo.
        \item Engenharia e implementação de um sistema de replay para arenas esportivas customizado para gravação de lances e transmissões de vídeo, hospedado em servidor VPS e gerenciado integralmente por uma API estruturada em Flask.
        \item Integração de fluxos de mídia de câmeras IP através dos protocolos RTSP e RTMP, permitindo a captura em tempo real, gravação contínua de replays estruturados e a inicialização de transmissões ao vivo.
        \item Automação de infraestruturas de segurança integrando visão computacional para o reconhecimento automático de placas de veículos (LPR).
        \item Projetos focados em Hardware e IoT, prototipando componentes eletrônicos e desenvolvendo firmware para dispositivos e microcontroladores Arduino.
    \end{itemize}

\cventry{Instituto Mauá de Tecnologia}{São Caetano do Sul, SP}{Trabalho de Conclusão de Curso - Dispositivo Wearable com Machine Learning}{Jan 2025 - Dez 2025}
    \begin{itemize}
        \item Desenvolvimento de sistema vestível preditivo com ESP32 para monitoramento de sinais fisiológicos (SpO2, frequência cardíaca) e fatores ambientais em tempo real.
        \item Implementação de modelos de Machine Learning (Random Forest, Árvore de Decisão) e análise acústica com TensorFlow (YAMNet) para predição de crises asmáticas.
        \item Criação de arquitetura de telemetria via protocolo MQTT e desenvolvimento de interface web utilizando Flask, HTML, CSS e JavaScript.
        \item Integração de assistente virtual com LLM (Google Gemini) baseado em arquitetura RAG para fornecer recomendações personalizadas ao paciente.
    \end{itemize}

% ------

\section{Habilidades}
    \begin{itemize}
        \item \textbf{Desenvolvimento Web e Software:} Python, TypeScript, JavaScript, React, SQL, HTML, CSS, Flask, Kivy, Supabase, Front-end, Back-end, FastAPI, APIs.
        \item \textbf{IA, Dados e Visão Computacional:} LLMs (OpenAI, Gemini, Ollama, Claude), Agentes de IA, RAG, Machine Learning, LPR, TensorFlow, Power BI.
        \item \textbf{Engenharia, IoT e Automação:} n8n, Arduino, ESP32, Arduino IDE, MQTT, Node-RED, SolidWorks, RPA, Git, Docker, ClickUp, CodeSys.
        \item \textbf{Idiomas:} Português (Nativo), Inglês (Avançado), Espanhol (Básico).
    \end{itemize}

\end{document}
